% Generated draft: requires mathematical and source review.
\documentclass[11pt]{article}
\usepackage[margin=25mm]{geometry}
\usepackage{fontspec}
\setmainfont{Latin Modern Roman}
\usepackage{amsmath,amssymb,mathtools,unicode-math}
\setmathfont{Latin Modern Math}
\usepackage{enumitem,needspace,xurl,hyperref}
\setlength{\parindent}{0pt}
\setlength{\parskip}{5pt}
\setlength{\emergencystretch}{3em}
\newcommand{\N}{\mathbb{N}}
\newcommand{\Z}{\mathbb{Z}}
\newcommand{\Q}{\mathbb{Q}}
\newcommand{\R}{\mathbb{R}}
\newcommand{\C}{\mathbb{C}}
\newcommand{\F}{\mathbb{F}}
\newcommand{\A}{\mathbb{A}}
\newcommand{\PP}{\mathbb{P}}
\newcommand{\E}{\mathbb{E}}
\newcommand{\eps}{\varepsilon}
\newcommand{\dd}{\,\mathrm{d}}
\begin{document}
\section*{UnsolvedMath ID: 11000204}
\textbf{Source problem code:} catalogue entry
\subsection{Definitions and mathematical statement}
Let $\mathcal M_{g,m}$ be the source's mapping class group of a genus-$g$ surface with $m$ marked points or punctures, with the precise marking and orientation conventions to be confirmed. A translation or half-translation structure determines a Teichm\"uller disk. Its orientation-preserving affine self-maps induce a subgroup of the mapping class group preserving that disk, and their derivative action gives the associated Veech group in $\mathrm{SL}_2(\R)$ or $\mathrm{PSL}_2(\R)$, subject to the usual kernel and sign conventions. In the common lattice-surface terminology, a Veech surface has finite-covolume Veech group.

\textbf{Problem 5.} Develop a theory of geometrically finite subgroups of $\mathcal M_{g,m}$ that includes the mapping-class subgroups defined by Veech surfaces. In particular, a proposed geometric-finiteness notion must explain how the affine/Teichm\"uller-disk subgroups arising from the intended Veech surfaces qualify. The requirement is not merely to study the derivative Fuchsian groups abstractly; the subgroups live in $\mathcal M_{g,m}$.

The excerpt gives no existing definition of geometric finiteness to characterize, nor a fixed theorem or list of axioms. The task is to formulate a useful class and its theory subject to the specified inclusion requirement. It must not be silently replaced by convex cocompactness or a purely pseudo-Anosov requirement: the intended Veech examples can have cusp/parabolic behavior and associated reducible mapping classes. Whether all affine stabilizers, lattice Veech groups only, or another standard Veech-surface convention is intended remains to be checked.
\subsection{Short English statement}
Develop an appropriate geometric-finiteness theory for mapping-class subgroups that includes the affine groups associated with Veech surfaces, including their possible cusp behavior.
\subsection{Sources}
\begin{itemize}
\item Local UnsolvedMath registry, permanent ID 11000204, source code catalogue entry; no proposer or proposal year supplied.
\item Farb, edited volume \emph{Problems on Mapping Class Groups and Related Topics}, chapter ``Geometric properties of the mapping class group,'' Problem 5, PDF page 247; \url{https://www.math.uchicago.edu/~farb/papers/mcgbook.pdf}. The registry records PDF-font-structure extraction and access on 2026-07-29; the inventory distinguishes 329 labeled headings by source page.
\item Veech surfaces are the source's named motivating examples. No separate primary reference or explicit proposed geometric-finiteness definition is supplied.
\end{itemize}
\end{document}
